\chapter{The Buffer Pool}
\label{ch:buffer}

The database is larger than memory, but queries must run at memory speed. The
\term{buffer pool} bridges the two: a fixed array of in-memory \term{frames},
each able to hold one page, plus the bookkeeping that decides which pages are
in memory and when they go back to disk.

\section{Frames, the page table, and pins}

\begin{figure}[ht]
\centering
\begin{tikzpicture}[font=\scriptsize]
\node[draw,minimum width=2.4cm,minimum height=2.4cm,align=left,fill=boxbg] (pt) {\textbf{page table}\\page 7 $\to$ frame 0\\page 3 $\to$ frame 2\\page 12 $\to$ frame 1};
\foreach \i/\p/\pin/\d in {0/7/2/clean,1/12/0/dirty,2/3/0/clean,3/--/--/free}{
  \node[draw,minimum width=2.2cm,minimum height=0.55cm,fill=white] (f\i) at (4.4,1.0-0.6*\i) {frame \i: page \p};
  \node[right=2pt of f\i,align=left,minimum width=1.9cm] (s\i) {pin \pin, \d};
}
\node[draw,cylinder,shape border rotate=90,aspect=0.2,minimum height=1.6cm,minimum width=1.4cm,fill=white] (disk) at (10.4,-0.1) {file};
\draw[arrow] (s1.east) -- node[above]{write back} (disk.west |- s1.east);
\end{tikzpicture}
\caption{A buffer pool with four frames. The page table maps page ids to frames; each frame has a pin count and a dirty bit.}
\label{fig:bpm}
\end{figure}

Each frame carries two pieces of state:

\begin{description}[leftmargin=1.2em]
\item[Pin count.] How many threads are using the page right now. A pinned page must not be evicted: someone holds a pointer into it.
\item[Dirty bit.] Whether the in-memory copy differs from the one on disk. A dirty page must be written before its frame is reused.
\end{description}

The interface follows from that:

\begin{itemize}
\item \code{FetchPage(id)}: if the page is in a frame, increment its pin and return it (a \term{hit}); otherwise choose a victim frame, write it back if dirty, read the page into it, pin it (a \term{miss}).
\item \code{UnpinPage(id, dirty)}: decrement the pin; mark the frame dirty if the caller changed it.
\item \code{NewPage()}: allocate a page id on disk and return a zeroed, pinned frame.
\item \code{FlushPage(id)}: write a dirty page now.
\end{itemize}

\begin{keyidea}
A pin is a promise, and a leaked pin is a leak of memory the pool can never
reclaim. Engines therefore wrap pins in a scoped object (a \term{page guard})
that unpins exactly once when it goes out of scope.
\end{keyidea}

\section{Choosing a victim}

When every frame is occupied and a miss arrives, the pool must evict an
unpinned page. The choice matters: a bad policy turns hits into misses, and
every miss is a disk read.

\paragraph{LRU.} Evict the least recently used page. Simple, and good for many
workloads. Its weakness is \term{sequential flooding}: one large scan touches
every page once, and each touch makes a cold scan page look ``recent'', so the
scan flushes out the hot index pages that every other query needs.

\paragraph{CLOCK.} An approximation of LRU that avoids moving pages in a list on
every access: each frame has a reference bit set on access, and a ``clock hand''
sweeps the frames, clearing bits and evicting the first frame whose bit is
already clear. Cheap, and still vulnerable to flooding.

\paragraph{LRU-K.} O'Neil, O'Neil, and Weikum's fix for flooding \cite{lruk}.
Remember the last $K$ access times of each page. The \term{backward
$K$-distance} of a page is the time since its $K$-th most recent access; a page
accessed fewer than $K$ times has infinite distance. Evict the page with the
largest distance.

With $K=2$, a page touched once by a scan has infinite distance, while an index
root touched thousands of times has a small one. The scan page goes first. LRU
is the special case $K = 1$.

\begin{table}[ht]
\centering
\begin{tabular}{@{}lccc@{}}
\toprule
Page & Accesses (time) & 2nd most recent & Backward 2-distance at $t=10$ \\
\midrule
A (index root) & 1, 4, 8, 9 & 8 & 2 \\
B (index leaf) & 3, 6 & 3 & 7 \\
C (scanned once) & 7 & --- & $\infty$ \\
\bottomrule
\end{tabular}
\caption{LRU-K with $K=2$. LRU would evict B (least recently used at time 6); LRU-K evicts C, the page with no second access, even though C was touched more recently than B.}
\label{tab:lruk}
\end{table}

When several pages have infinite distance, a tie-breaking rule is needed.
A common choice is the page whose oldest remembered access is earliest; with
$K = 2$ such a page has exactly one access, so this is plain LRU among the
pages seen once.

\section{Steal and force}

Two policy questions connect the buffer pool to recovery, and their answers
decide what the log must contain.

\begin{description}[leftmargin=1.2em]
\item[Steal?] May the pool write a dirty page of an \emph{uncommitted} transaction to disk, to free its frame? If yes, the disk can hold changes that must later be \emph{undone}.
\item[Force?] Must every page a transaction changed be written to disk \emph{before} it commits? If no, a committed change may exist only in memory at the time of a crash, and must later be \emph{redone}.
\end{description}

\begin{table}[ht]
\centering
\begin{tabular}{@{}lll@{}}
\toprule
& Force & No force \\
\midrule
No steal & no undo, no redo (slow, memory-bound) & redo only \\
Steal & undo only & undo and redo (fast, flexible) \\
\bottomrule
\end{tabular}
\caption{The four combinations. High-performance engines choose steal and no force, and pay for it with a log that supports both undo and redo.}
\label{tab:stealforce}
\end{table}

Steal lets the pool evict any unpinned page, so a long transaction cannot
exhaust memory. No force lets a commit write one small log record instead of
every page it touched. Both are wins, and \cref{ch:wal,ch:aries} show how the
log makes them safe.

\begin{ontodb}[: the buffer pool]
OntoDB's pool uses \textbf{LRU-K} with $K = 2$ by default (\code{\textbackslash set lru\_k}),
tie-breaking by the older $K$-th access. Its policy is \textbf{steal / no force}.
Pages are reached only through \code{ReadPageGuard} (shared latch, one pin)
and \code{WritePageGuard} (exclusive latch, one pin); both unpin exactly once.
The test harness asserts that the number of pinned frames is zero after every
query and every destroyed iterator: a leaked pin fails the test even when the
query's rows were right.

When the pool writes a dirty page and a log manager is installed, the order is
fixed: first flush the log up to the page's pageLSN, then write the page. A
named crash point sits between the two, so a test can kill the process after the
log is durable and before the page is. That ordering is the write-ahead rule of
\cref{ch:wal}.
\end{ontodb}

\begin{decision}{One I/O thread, first in, first out}
\dpart{Choice} The pool never calls the disk manager itself. It hands a request
to a \code{DiskScheduler}, whose single background thread runs reads and writes
in the order they were scheduled. \code{Schedule} returns at once with a
future; if the I/O fails, the error is rethrown from the future's \code{get}.
\dpart{Alternatives} Synchronous I/O on the query thread; a pool of I/O
threads; an asynchronous kernel interface such as \code{io\_uring}.
\dpart{Why} In-order execution makes ``write page 9, then read page 9'' correct
with no per-page bookkeeping. A background thread lets a flush or a prefetch
overlap with query work, and errors arrive where the waiting caller can handle
them.
\dpart{Cost} No parallel I/O: an SSD reaches full throughput only with many
requests in flight, and one slow read delays every request behind it.
\end{decision}

\begin{decision}{LRU-K with $K = 2$}
\dpart{Choice} Evict the unpinned frame with the largest backward 2-distance.
Frames seen fewer than twice have infinite distance and go first, earliest
oldest access first. Time is a logical clock that ticks on every access, not
the wall clock, and a pinned frame keeps its history.
\dpart{Alternatives} LRU, CLOCK, 2Q, ARC.
\dpart{Why} SPARQL plans scan: an unselective pattern reads every leaf of a
permutation once. LRU and CLOCK let that scan push out the index's upper levels;
LRU-K does not. A logical clock and a fixed tie-break make eviction
deterministic, so a test can predict exactly which frame goes; the replacer's
tests replay \cref{tab:lruk} and check that C, not B, is evicted. Keeping
history while a frame is pinned stops a hot page from looking new each time it
is unpinned.
\dpart{Cost} Each frame keeps $K$ timestamps, and choosing a victim scans every
frame: linear in the pool size, which is cheap for the default 64 frames but
would need a priority structure for pools of millions. A second policy is
planned as a measured comparison, not as a replacement.
\end{decision}

\begin{decision}{One latch for the whole pool, held across I/O}
\dpart{Choice} A single mutex covers the page table, the free list, every pin
count and dirty bit, and the replacer calls. It stays held while a miss reads
its page or a victim is written back.
\dpart{Alternatives} Partition the page table and give each frame an
``I/O in progress'' state, so other threads keep working during a read; this is
roughly what PostgreSQL does with buffer-mapping partitions and per-buffer I/O
locks.
\dpart{Why} The pool's state is then always consistent: no frame is
half-loaded while another thread looks it up, and the tests pass under
ThreadSanitizer without subtle protocols. With a single I/O thread
(design decision on the scheduler), finer locking would mostly move the
waiting elsewhere.
\dpart{Cost} A miss stalls every thread for one disk read. Finer-grained
latching is the first thing to change when I/O becomes parallel.
\end{decision}

\begin{decision}{The pool's destructor does not flush}
\dpart{Choice} Destroying the pool discards dirty frames. Durable shutdown
calls \code{FlushAllPages} explicitly (plan 3.6).
\dpart{Alternatives} Flush every dirty page on destruction, which is
convenient.
\dpart{Why} Durability is the log's job. Tests need a path that behaves like a
dead process, writing nothing, and a destructor also runs while an exception
from a crash point unwinds; a flushing destructor would write exactly the page
the test is trying to keep off disk.
\dpart{Cost} A caller who forgets \code{FlushAllPages} loses work on a clean
exit. The clean-restart test of plan 3.6 is there to catch it.
\end{decision}

\begin{decision}{A full pool returns null}
\dpart{Choice} When every frame is pinned, \code{NewPage} and
\code{FetchPage} return \code{nullptr} at once.
\dpart{Alternatives} Block until some frame is unpinned, or throw.
\dpart{Why} The caller usually holds pins of its own; blocking could wait for
a pin that only the waiter can release. Returning lets the caller give up the
operation and release what it holds.
\dpart{Cost} Every caller must check for null; guards and the tree turn that
into an error close to the cause.
\end{decision}

\begin{decision}{Pages only through guards}
\dpart{Choice} Code above the pool reaches a page only through a
\code{ReadPageGuard} or \code{WritePageGuard}: one pin plus a latch, released
exactly once by \code{Drop} or the destructor. Guards move but never copy, and a
moved-from guard is empty. Tests assert the pinned-frame count is zero after
every query.
\dpart{Alternatives} Callers call \code{UnpinPage} by hand, as early textbook
pools do.
\dpart{Why} A leaked pin returns the right rows and silently shrinks the pool;
a double unpin lets a page in use be evicted. Neither is visible in query
results. Tying the pin to a scope removes both.
\dpart{Cost} Ownership rules (move-only, drop before pinning the sibling) that
every caller must follow, and a little ceremony for one-line accesses.
\end{decision}

\begin{decision}{Pin, then latch; unlatch, then unpin}
\dpart{Choice} Each frame has a reader--writer latch. A guard asks the pool to
pin the page, then takes the latch, outside the pool's own latch. Dropping
releases the latch first and the pin second. A null page from the pool becomes
an error at the call that asked for it.
\dpart{Alternatives} Take the page latch while still holding the pool latch;
or release the pin before the latch.
\dpart{Why} Waiting for a page latch while holding the pool latch would block
every thread behind one busy page, and could deadlock with the thread that
holds the page and needs the pool. Unpinning first would let the frame be
evicted and reused while its latch is still held.
\dpart{Cost} Between the pin and the latch, another writer may get the page
first; callers must not assume what they will find.
\end{decision}

\begin{decision}{A latch belongs to a page, not to a frame}
\dpart{Choice} When a frame is given to a different page, it gets a new latch.
This happens only while the frame is unpinned, and since guards pin before they
latch, no thread holds or waits on the old one.
\dpart{Alternatives} One latch per frame for the life of the pool.
\dpart{Why} Latch order is a rule about pages: top-down and left to right. With
a latch per frame, two unrelated pages that once shared a frame look like the
same lock, and ThreadSanitizer reported a lock-order cycle in a bulk load on a
three-frame pool that could never deadlock. A fresh latch makes the checker see
the same thing the protocol is about.
\dpart{Cost} A small allocation on every page load, next to a disk read that
costs far more. (It has to be a new allocation: ThreadSanitizer forgets a mutex
when its memory is freed, but not when a statically initialized
\code{std::shared\_mutex} is destroyed and rebuilt in place, which was the first
attempt.)
\end{decision}

\begin{decision}{A write guard dirties only if written}
\dpart{Choice} A write guard marks the page dirty on drop only if its mutable
\code{GetData} was called.
\dpart{Alternatives} Every write guard marks its page dirty.
\dpart{Why} B+ tree latch crabbing takes exclusive latches on pages it may
never change. Marking those dirty would write clean pages back to disk.
\dpart{Cost} Code that writes through a pointer obtained earlier, after
dropping the guard, escapes the dirty bit, which is one more reason never to
keep that pointer.
\end{decision}

\begin{decision}{Inspecting a page goes through the pool}
\dpart{Choice} \code{\textbackslash page N} reads page $N$ with a read guard,
as any operator would, and prints its pageLSN and the first payload bytes.
\code{\textbackslash bpm} prints each frame's page, pins, dirty bit, and
pageLSN.
\dpart{Alternatives} Read the page straight from the file, bypassing the pool.
\dpart{Why} The pool's copy is what the engine sees, including changes that
are dirty and not yet on disk. Reading the file would show a page the engine
has already moved past.
\dpart{Cost} Looking is an access: it can load the page, evict another, and add
to the page's LRU-K history. The observer changes what it observes, so
\code{\textbackslash bpm} is best read before \code{\textbackslash page}.
\end{decision}

\begin{status}
The disk manager, the I/O scheduler, the LRU-K replacer, and the pool run
(plans 1.1 to 1.4), and so do page guards with per-frame reader--writer latches
(plan 1.5) and the shell's \code{\textbackslash bpm}, \code{\textbackslash page},
and \code{\textbackslash trace} (plan 1.6): the buffer pool is complete. It
already flushes the log up to a page's pageLSN before writing the page. The
default memory backend has no pool, so these commands come alive with the
indexed store (plan 3.4).
The experiment that checks them is simple: run a scan with a pool of 4 frames
and of 64, with $K=1$ and $K=2$, and compare hits and misses in \code{\textbackslash stats}.
\end{status}

\begin{checkyourself}
\item A frame has pin count 0 and is dirty. What must happen before a different page is read into it?
\item In \cref{tab:lruk}, add an access to C at time 10. Which page does LRU-K evict at time 11 now?
\item With $K = 3$, frame X has accesses at times 2 and 9, and frame Y only at time 5. Both are unpinned. Which does OntoDB evict, and why is X's more recent access irrelevant?
\item Explain sequential flooding, and why $K = 2$ resists it while $K = 1$ does not.
\item Which cell of the steal/force table needs \emph{undo} information in the log, and why?
\item A thread schedules a write of page 9 and then, without waiting, a read of page 9 into another buffer. What does the read return in OntoDB, and what could go wrong with a pool of I/O threads?
\item Two threads miss on two different pages at the same moment. In OntoDB's pool, how long does the second thread wait, and why?
\item A test arms the crash point between the log flush and the page write, and the exception unwinds through the pool's destructor. Why would a destructor that flushes dirty pages ruin the test?
\item Why does a guard release its latch \emph{before} it unpins the page, and not after?
\item You run \code{\textbackslash page 7} twice and then \code{\textbackslash bpm}. How has inspecting the page changed the pool?
\item A query returns the right rows, but afterwards one frame is still pinned. What will go wrong later, and how do OntoDB's guards prevent it?
\end{checkyourself}

\section*{Further reading}
Petrov \cite{petrov}, chapter 5 on buffer management. O'Neil et al.\ \cite{lruk}.
Effelsberg and H\"arder \cite{effelsberg} on buffer management principles.
